\section{Vocabulaire}

\begin{exemplebox}
\includegraphics[
    width=\linewidth,
    keepaspectratio
]{images/droite_graduée.png}%
\end{exemplebox}
\begin{definitionbox}
\begin{itemize}
    \item L'ensemble des nombres \hlb{relatifs} est la réunion de l'ensemble
        des nombres \answer{positifs} et de l'ensemble des nombres
        \answer{négatifs}.
    \item Sur une droite graduée, tout point est repéré par un nombre relatif
        appelé \answer{abscisse}. 
    \item Pour un nombre donné, sa \answer{distance à zéro} est la distance
        entre le point d'abscisse 0 et le point d'abscisse le nombre donné.
    \item La distance à zéro d'un nombre est également appelée \answer{sa valeur
        absolue}.
    \item Deux \answer{nombres opposés} ont la même distance à zéro et sont de
        signes contraire.
\end{itemize}
\end{definitionbox}
\begin{exempleboxsuite}
    \begin{itemize}
        \item A a pour \answer{abscisse} 3, on le note A(3) 
        \item (-4) et (+4) ont la même \answer{distance à zéro}. 
        \item La \answer{valeur absolue} de (-4) est 4. 
        \item On dit que (-4) et (+4) sont des \answer{nombres opposés}.
    \end{itemize}
\end{exempleboxsuite}




